% !TeX root = ../../Book3.tex
% 纲要标题：### 6.2 弱零点定理
% 纲要位置：工作记录/计划纲要.md，第 2070 行。
% 纲要细目（写作范围）：
% * 代数闭域上的极大理想；多变量求值映射的核
% * 公共零点的存在判据；坐标点与素谱点的区别
% * 非代数闭域的对照

\section{弱零点定理}\label{b3:sec:0602}

一个坐标点可以由所有在该点取零的多项式描述。设 \(k\) 为任意域，\(n\ge0\)，\(a=(a_1,\ldots,a_n)\in k^n\)，考虑取值映射 \(\operatorname{ev}_a:k[x_1,\ldots,x_n]\to k\)。逐次将 \(x_1,\ldots,x_n\) 替换成 \(a_1,\ldots,a_n\)，每次作差都是关于 \(x_i\) 的多项式在 \(x_i\) 与 \(a_i\) 处的取值之差，因而被 \(x_i-a_i\) 整除。把这些差相加，对任意 \(f\) 得到
\[
f(x_1,\ldots,x_n)-f(a_1,\ldots,a_n)
=\sum_{i=1}^n(x_i-a_i)g_i(x_1,\ldots,x_n),
\qquad g_i\in k[x_1,\ldots,x_n].
\]
因此 \(\ker\operatorname{ev}_a=(x_1-a_1,\ldots,x_n-a_n)\)。取值映射在常数上是恒等映射，因而满射到域 \(k\)，所以这个核是极大理想。代数闭域上的弱零点定理将说明，每个极大理想都以这种方式产生。

\subsection{代数闭域上的极大理想}\label{b3:subsec:0602:01}

\begin{theorem}[Hilbert 弱零点定理]\label{b3:thm:weak-nullstellensatz}
设 \(k\) 是代数闭域，\(n\ge0\)。则 \(k[x_1,\ldots,x_n]\) 的每个极大理想唯一写成
\[
\mathfrak m_a=(x_1-a_1,\ldots,x_n-a_n),\qquad a\in k^n.
\]
这称为\term[ruo-lingdian-dingli]{弱零点定理}[weak Nullstellensatz]。
\end{theorem}
\begin{proof}
对极大理想 \(\mathfrak m\)，由\cref{b3:cor:finite-type-residue-fields}及代数闭性，结构映射 \(k\to k[x_1,\ldots,x_n]/\mathfrak m\) 是同构。令 \(a_i\in k\) 为 \(x_i\) 在商域中的像，则 \(x_i-a_i\in\mathfrak m\)，故 \(\mathfrak m_a\subseteq\mathfrak m\)。二者都极大，只能相等。若 \(\mathfrak m_a=\mathfrak m_b\)，将 \(x_i-a_i\) 在 \(b\) 处取值，得到 \(b_i=a_i\)，所以点唯一。
\end{proof}

\subsection{公共零点的存在判据}\label{b3:subsec:0602:02}

对理想 \(I\subseteq k[x_1,\ldots,x_n]\)，记
\[
Z_k(I)=\bigl\{\,a\in k^n\mid f(a)=0\text{ 对每个 }f\in I\,\bigr\}.
\]
\symidx{\(Z_k(I)\)：理想 \(I\) 在 \(k^n\) 中的公共零点集}
这与素谱中的 \(V(I)\) 需要区别：\(Z_k(I)\) 只记录 \(k\) 坐标点，而 \(V(I)\) 包含所有含 \(I\) 的素理想。若 \(k\) 代数闭，弱零点定理把 \(Z_k(I)\) 中的坐标点与 \(V(I)\) 中的闭点对应起来。例如取 \(I=(0)\subset k[x,y]\)，则 \(Z_k(0)=k^2\)，而 \(V(0)=\operatorname{Spec}k[x,y]\) 还包含 \((0),(x),(y-x^2)\) 等非闭点；它们的商环分别为 \(k[x,y],k[y],k[x]\)，都是整环却不是域，所以这些素理想都不极大。这些非闭点即使在 \(k\) 代数闭时也仍然存在。

\begin{corollary}\label{b3:cor:proper-ideal-common-zero}
设 \(k\) 为代数闭域，\(n\ge0\)，\(I\) 为 \(k[x_1,\ldots,x_n]\) 的理想。记 \(Z_k(I)=\{a\in k^n:f(a)=0\text{ 对所有 }f\in I\}\)。则 \(I\) 是真理想，当且仅当 \(Z_k(I)\ne\varnothing\)。特别地，给定整数 \(r\ge1\) 及 \(k[x_1,\ldots,x_n]\) 中的多项式 \(f_1,\ldots,f_r\)，它们没有 \(k^n\) 中的公共零点，当且仅当存在 \(g_1,\ldots,g_r\in k[x_1,\ldots,x_n]\) 满足
\[
g_1f_1+\cdots+g_rf_r=1.
\]
\end{corollary}
\begin{proof}
若 \(I\) 是真理想，可先取极大理想 \(\mathfrak m\supseteq I\)。\cref{b3:thm:weak-nullstellensatz}把这个纯代数的选择转为一个坐标点：\(\mathfrak m=\mathfrak m_a\)，其中 \(a\in k^n\)。于是每个 \(f\in I\) 都属于 \(\ker\operatorname{ev}_a\)，即 \(f(a)=0\)，所以 \(a\) 是公共零点。反过来，若 \(a\) 为公共零点，则 \(I\subseteq\ker\operatorname{ev}_a\)，所以 \(1\notin I\)。最后对 \(I=(f_1,\ldots,f_r)\) 使用理想生成的定义即可。
\end{proof}

例如 \(xy-1\) 和 \(x\) 没有公共零点，等式
\(yx-(xy-1)=1\) 就是可以直接核验的证据。零点定理保证任何无公共零点的有限系统都有这样的代数组合；它本身没有给出这些系数的次数上界。第一卷附录F.3的 Buchberger 算法及定理F.3.3给出一种计算办法：系数域支持精确四则运算和零判定，且选定的全局单项式序能够有效比较时，可从有限生成元计算 Gröbner 基；保存各步的生成元表示，再把 \(1\) 的约化过程回代，就能得到上述理想成员证据。

\subsection{非代数闭域上的对照}\label{b3:subsec:0602:03}

\((x^2+1)\subseteq\mathbb R[x]\) 是真理想，却没有实零点。有限域上的差别甚至会出现在零理想的零点理想中：多项式 \(x^q-x\) 在全部 \(\mathbb F_q\) 元素上取零，本身却不是零多项式。这些例子分别说明公共零点存在性及从点集恢复多项式关系，都依赖所使用的基域。

\ProseParagraphBreak
不过任何域上的真理想，在代数闭包中都有公共零点。对真理想 \(I\subseteq k[x_1,\ldots,x_n]\)，选择极大理想 \(\mathfrak m\supseteq I\)，其剩余域 \(L\) 由\cref{b3:cor:finite-type-residue-fields}是有限代数扩张。把 \(L\) 嵌入一个固定代数闭包 \(\overline k\)，变量的像给出所求点。这里有限扩域嵌入代数闭包，可按有限组代数生成元逐次选择最小多项式的根完成；不要求 \(L/k\) 可分。

\ProseParagraphBreak
若只研究实点或有限域上的有理点，就应当保留相应基域的额外关系，而不能把代数闭域版本原样照用。本卷后续的实代数补充也将据此区分不同的零点定理。
